\documentclass[11pt,a4paper]{article}

% Packages
\usepackage[utf8]{inputenc}
\usepackage[T1]{fontenc}
\usepackage[margin=2.5cm]{geometry}
\usepackage{graphicx}
\usepackage{xcolor}
\usepackage{listings}
\usepackage{hyperref}
\usepackage{fancyhdr}
\usepackage{tcolorbox}
\usepackage{enumitem}
\usepackage{tikz}

% TikZ libraries
\usetikzlibrary{shapes,arrows,positioning}

% Code listing style
\lstset{
    language=C,
    basicstyle=\ttfamily\small,
    keywordstyle=\color{blue}\bfseries,
    commentstyle=\color{green!60!black}\itshape,
    stringstyle=\color{red},
    showstringspaces=false,
    breaklines=true,
    frame=single,
    numbers=left,
    numberstyle=\tiny\color{gray},
    tabsize=2
}

% Color boxes
\tcbuselibrary{skins,breakable}
\newtcolorbox{objectivebox}{
    colback=blue!5!white,
    colframe=blue!75!black,
    title=Lab Objectives,
    fonttitle=\bfseries
}

\newtcolorbox{prereqbox}{
    colback=yellow!5!white,
    colframe=yellow!75!black,
    title=Prerequisites,
    fonttitle=\bfseries
}

\newtcolorbox{solutionbox}{
    colback=green!5!white,
    colframe=green!75!black,
    title=Solution,
    fonttitle=\bfseries,
    breakable
}

\newtcolorbox{warningbox}{
    colback=red!5!white,
    colframe=red!75!black,
    title=Warning,
    fonttitle=\bfseries
}

% Header and footer
\pagestyle{fancy}
\fancyhf{}
\lhead{USP Training Course 1}
\rhead{Lab 1.1: Basic LoRaWAN Join}
\cfoot{\thepage}

\title{
    \textbf{Lab 1.1: Basic LoRaWAN Join}\\
    \large Course 1: LoRa \& LoRaWAN Fundamentals
}
\author{Dr. ABDELMALEK OMAR}
\date{}

\begin{document}

\maketitle
\thispagestyle{fancy}

% ============================================================================
\section{Lab Overview}
% ============================================================================

\begin{objectivebox}
By the end of this lab, you will be able to:
\begin{itemize}
    \item Configure LoRaWAN credentials (DevEUI, JoinEUI, AppKey)
    \item Perform OTAA (Over-The-Air Activation) join procedure
    \item Understand the join accept message
    \item Monitor join status through serial console
    \item Verify device registration on network server
\end{itemize}
\end{objectivebox}

\vspace{0.5cm}

\begin{prereqbox}
\textbf{Required Knowledge:}
\begin{itemize}
    \item Basic understanding of LoRaWAN architecture
    \item Familiarity with Zephyr RTOS build system
    \item Access to LoRaWAN network server (TTN, ChirpStack, etc.)
\end{itemize}

\vspace{0.3cm}

\textbf{Hardware Required:}
\begin{itemize}
    \item Xiao nRF54L15 development board
    \item LR1120 LoRa shield
    \item USB cable for programming and serial console
    \item LoRaWAN gateway within range
\end{itemize}

\vspace{0.3cm}

\textbf{Software Required:}
\begin{itemize}
    \item Zephyr SDK installed
    \item West tool configured
    \item Serial terminal (screen, minicom, or PuTTY)
\end{itemize}
\end{prereqbox}

% ============================================================================
\section{Lab Duration}
% ============================================================================

\textbf{Estimated Time:} 45 minutes

\begin{itemize}
    \item Setup and configuration: 15 minutes
    \item Implementation: 20 minutes
    \item Testing and verification: 10 minutes
\end{itemize}

% ============================================================================
\section{Background}
% ============================================================================

\subsection{OTAA Join Procedure}

LoRaWAN supports two activation methods:
\begin{itemize}
    \item \textbf{OTAA (Over-The-Air Activation):} Device performs join procedure to obtain session keys
    \item \textbf{ABP (Activation By Personalization):} Session keys pre-configured in device
\end{itemize}

OTAA is the recommended method for production deployments because:
\begin{itemize}
    \item More secure (keys derived from join procedure)
    \item Supports key rotation
    \item Enables device mobility between networks
\end{itemize}

\subsection{Join Request Message}

The join request contains:
\begin{itemize}
    \item \textbf{JoinEUI} (formerly AppEUI): 64-bit identifier for join server
    \item \textbf{DevEUI}: 64-bit globally unique device identifier
    \item \textbf{DevNonce}: 16-bit random value (prevents replay attacks)
\end{itemize}

The message is encrypted with the AppKey.

\subsection{Join Accept Message}

Upon successful join, the network server responds with:
\begin{itemize}
    \item \textbf{JoinNonce}: Random value from network
    \item \textbf{NetID}: Network identifier
    \item \textbf{DevAddr}: Device address (for this session)
    \item \textbf{DLSettings}: Downlink settings (RX1 delay, RX2 data rate)
    \item \textbf{RXDelay}: RX window timing
\end{itemize}

From this, the device derives:
\begin{itemize}
    \item \textbf{NwkSKey}: Network session key (for MAC commands)
    \item \textbf{AppSKey}: Application session key (for payload encryption)
\end{itemize}

% ============================================================================
\section{Lab Instructions}
% ============================================================================

\subsection{Step 1: Register Device on Network Server}

\begin{enumerate}
    \item Log in to your LoRaWAN network server (e.g., The Things Network)
    \item Create a new application (if not already exists)
    \item Add a new end device:
    \begin{itemize}
        \item Activation mode: OTAA
        \item LoRaWAN version: 1.0.4
        \item Regional Parameters: EU868 (or your region)
        \item DevEUI: Copy from device label or generate
        \item JoinEUI: Use your application's JoinEUI
        \item AppKey: Generate a new 128-bit key
    \end{itemize}
    \item Save the credentials for the next step
\end{enumerate}

\subsection{Step 2: Configure Device Credentials}

Navigate to your project directory and edit the configuration file:

\begin{lstlisting}[language=bash]
cd samples/lora/lorawan_hello
\end{lstlisting}

Create or edit \texttt{prj.conf}:

\begin{lstlisting}[language=bash]
# LoRaWAN Configuration
CONFIG_LORAWAN=y
CONFIG_LORAWAN_REGION_EU868=y
CONFIG_LORAWAN_ACTIVATION_OTAA=y

# Enable logging
CONFIG_LOG=y
CONFIG_LORAWAN_LOG_LEVEL_DBG=y
\end{lstlisting}

Edit \texttt{src/main.c} to add your credentials (see Solution section).

\subsection{Step 3: Build the Firmware}

\begin{lstlisting}[language=bash]
west build -b xiao_nrf54l15 samples/lora/lorawan_hello
\end{lstlisting}

\subsection{Step 4: Flash the Device}

\begin{lstlisting}[language=bash]
west flash
\end{lstlisting}

\subsection{Step 5: Monitor Serial Output}

Open a serial terminal:

\begin{lstlisting}[language=bash]
# Linux/macOS
screen /dev/ttyUSB0 115200

# Or use minicom
minicom -D /dev/ttyUSB0 -b 115200
\end{lstlisting}

\subsection{Step 6: Observe Join Procedure}

You should see output similar to:

\begin{lstlisting}
[00:00:00.100,000] <inf> lorawan: LoRaWAN stack initialized
[00:00:00.150,000] <inf> lorawan: Starting OTAA join procedure
[00:00:05.230,000] <inf> lorawan: Join request sent
[00:00:10.450,000] <inf> lorawan: Join accept received
[00:00:10.451,000] <inf> lorawan: Network joined successfully!
[00:00:10.452,000] <inf> lorawan: DevAddr: 0x260B1234
\end{lstlisting}

\subsection{Step 7: Verify on Network Server}

\begin{enumerate}
    \item Return to network server console
    \item Navigate to your device page
    \item Check "Live Data" or "Events" tab
    \item Verify join request was received
    \item Confirm join accept was sent
    \item Device should show as "Active" or "Online"
\end{enumerate}

% ============================================================================
\section{Solution}
% ============================================================================

\begin{solutionbox}

\subsection{Complete main.c Implementation}

\begin{lstlisting}
#include <zephyr/kernel.h>
#include <zephyr/device.h>
#include <zephyr/logging/log.h>
#include <zephyr/lorawan/lorawan.h>

LOG_MODULE_REGISTER(lab_1_1, LOG_LEVEL_INF);

/* LoRaWAN Credentials - REPLACE WITH YOUR OWN */
#define LORAWAN_DEV_EUI    { 0x00, 0x00, 0x00, 0x00, \
                             0x00, 0x00, 0x00, 0x01 }
#define LORAWAN_JOIN_EUI   { 0x00, 0x00, 0x00, 0x00, \
                             0x00, 0x00, 0x00, 0x00 }
#define LORAWAN_APP_KEY    { 0x00, 0x00, 0x00, 0x00, \
                             0x00, 0x00, 0x00, 0x00, \
                             0x00, 0x00, 0x00, 0x00, \
                             0x00, 0x00, 0x00, 0x00 }

static uint8_t dev_eui[] = LORAWAN_DEV_EUI;
static uint8_t join_eui[] = LORAWAN_JOIN_EUI;
static uint8_t app_key[] = LORAWAN_APP_KEY;

static struct lorawan_join_config join_cfg = {
    .mode = LORAWAN_ACT_OTAA,
    .dev_eui = dev_eui,
    .otaa = {
        .join_eui = join_eui,
        .app_key = app_key,
        .nwk_key = app_key,  // For LoRaWAN 1.0.x
    }
};

/* LoRaWAN downlink callback */
static void dl_callback(uint8_t port, bool data_pending,
                       int16_t rssi, int8_t snr,
                       uint8_t len, const uint8_t *data)
{
    LOG_INF("Downlink received:");
    LOG_INF("  Port: %d", port);
    LOG_INF("  RSSI: %d dBm", rssi);
    LOG_INF("  SNR: %d dB", snr);
    LOG_INF("  Length: %d bytes", len);

    if (len > 0) {
        LOG_HEXDUMP_INF(data, len, "Payload");
    }
}

/* LoRaWAN datarate changed callback */
static void dr_changed_callback(enum lorawan_datarate dr)
{
    LOG_INF("Datarate changed to: DR%d", dr);
}

int main(void)
{
    int ret;
    const struct device *lora_dev;

    LOG_INF("==============================================");
    LOG_INF("Lab 1.1: Basic LoRaWAN Join");
    LOG_INF("==============================================");

    /* Get LoRa device */
    lora_dev = DEVICE_DT_GET(DT_ALIAS(lora0));
    if (!device_is_ready(lora_dev)) {
        LOG_ERR("LoRa device not ready");
        return -1;
    }

    LOG_INF("LoRa device ready");

    /* Initialize LoRaWAN stack */
    ret = lorawan_start();
    if (ret < 0) {
        LOG_ERR("Failed to start LoRaWAN stack: %d", ret);
        return ret;
    }

    LOG_INF("LoRaWAN stack started");

    /* Register callbacks */
    struct lorawan_downlink_cb downlink_cb = {
        .port = LW_RECV_PORT_ANY,
        .cb = dl_callback
    };
    lorawan_register_downlink_callback(&downlink_cb);
    lorawan_register_dr_changed_callback(dr_changed_callback);

    LOG_INF("Callbacks registered");

    /* Perform OTAA join */
    LOG_INF("Starting OTAA join procedure...");
    ret = lorawan_join(&join_cfg);
    if (ret < 0) {
        LOG_ERR("Join failed: %d", ret);
        return ret;
    }

    LOG_INF("Join request sent, waiting for accept...");

    /* Wait for join to complete */
    while (!lorawan_is_joined()) {
        k_sleep(K_SECONDS(1));
    }

    LOG_INF("==============================================");
    LOG_INF("SUCCESS: Network joined!");
    LOG_INF("==============================================");

    /* Get and display DevAddr */
    uint32_t dev_addr;
    lorawan_get_dev_addr(&dev_addr);
    LOG_INF("Assigned DevAddr: 0x%08X", dev_addr);

    /* Keep application running */
    while (1) {
        k_sleep(K_SECONDS(60));
        LOG_INF("Device active, joined to network");
    }

    return 0;
}
\end{lstlisting}

\subsection{Build Configuration (prj.conf)}

\begin{lstlisting}[language=bash]
# LoRaWAN Configuration
CONFIG_LORAWAN=y
CONFIG_LORAWAN_REGION_EU868=y
CONFIG_LORAWAN_ACTIVATION_OTAA=y

# Enable logging
CONFIG_LOG=y
CONFIG_LORAWAN_LOG_LEVEL_DBG=y
CONFIG_LOG_BUFFER_SIZE=2048

# Network stack
CONFIG_NETWORKING=y
CONFIG_NET_L2_LORAWAN=y

# Increase stack sizes
CONFIG_MAIN_STACK_SIZE=2048
CONFIG_SYSTEM_WORKQUEUE_STACK_SIZE=2048
\end{lstlisting}

\end{solutionbox}

% ============================================================================
\section{Expected Output}
% ============================================================================

When the device successfully joins, you should see:

\begin{lstlisting}
==============================================
Lab 1.1: Basic LoRaWAN Join
==============================================
[00:00:00.100] <inf> lab_1_1: LoRa device ready
[00:00:00.105] <inf> lab_1_1: LoRaWAN stack started
[00:00:00.110] <inf> lab_1_1: Callbacks registered
[00:00:00.115] <inf> lab_1_1: Starting OTAA join procedure...
[00:00:00.120] <inf> lab_1_1: Join request sent, waiting...
[00:00:05.234] <inf> lorawan: TX: Join Request
[00:00:05.235] <inf> lorawan: Frequency: 868.1 MHz
[00:00:05.236] <inf> lorawan: DR: SF12BW125
[00:00:10.567] <inf> lorawan: RX: Join Accept
[00:00:10.568] <inf> lorawan: Network joined
==============================================
SUCCESS: Network joined!
==============================================
[00:00:10.570] <inf> lab_1_1: Assigned DevAddr: 0x260B1234
[00:01:10.570] <inf> lab_1_1: Device active, joined to network
\end{lstlisting}

% ============================================================================
\section{Troubleshooting}
% ============================================================================

\subsection{Join Request Not Received by Gateway}

\begin{warningbox}
\textbf{Symptoms:} No join request appears on network server

\textbf{Possible Causes:}
\begin{itemize}
    \item Gateway out of range
    \item Incorrect frequency plan (EU868 vs US915)
    \item Antenna not connected
    \item Radio not configured properly
\end{itemize}

\textbf{Solutions:}
\begin{itemize}
    \item Check gateway status and location
    \item Verify regional parameters match gateway
    \item Ensure antenna is securely connected
    \item Check device tree configuration for radio
\end{itemize}
\end{warningbox}

\subsection{Join Accept Not Received}

\begin{warningbox}
\textbf{Symptoms:} Join request sent but no accept received

\textbf{Possible Causes:}
\begin{itemize}
    \item Incorrect credentials (DevEUI, JoinEUI, AppKey)
    \item RX window timing issues
    \item Device not registered on network server
\end{itemize}

\textbf{Solutions:}
\begin{itemize}
    \item Double-check all credentials match network server
    \item Verify device is registered and enabled
    \item Check RX window configuration
    \item Ensure AppKey is correct (case-sensitive)
\end{itemize}
\end{warningbox}

\subsection{Build Errors}

\begin{warningbox}
\textbf{Common Build Issues:}
\begin{itemize}
    \item Missing \texttt{CONFIG\_LORAWAN=y} in prj.conf
    \item Incorrect board name
    \item Missing Zephyr modules
\end{itemize}

\textbf{Solutions:}
\begin{itemize}
    \item Clean build: \texttt{west build -t pristine}
    \item Update Zephyr: \texttt{west update}
    \item Check board support: \texttt{west boards | grep xiao}
\end{itemize}
\end{warningbox}

% ============================================================================
\section{Deliverables}
% ============================================================================

Upon completion of this lab, you should submit:

\begin{enumerate}[itemsep=10pt]
    \item \textbf{Source Code}
    \begin{itemize}
        \item Complete \texttt{main.c} with your credentials
        \item \texttt{prj.conf} configuration file
    \end{itemize}

    \item \textbf{Serial Console Log}
    \begin{itemize}
        \item Full output showing successful join
        \item Include timestamps
        \item Highlight DevAddr assigned
    \end{itemize}

    \item \textbf{Network Server Screenshot}
    \begin{itemize}
        \item Device showing as "Active" or "Online"
        \item Join request/accept visible in events
        \item DevAddr matches console output
    \end{itemize}

    \item \textbf{Lab Report}
    \begin{itemize}
        \item Brief description of procedure followed
        \item Any issues encountered and how resolved
        \item Time taken for join procedure
        \item RSSI and SNR values observed
    \end{itemize}
\end{enumerate}

% ============================================================================
\section{Additional Challenges}
% ============================================================================

\subsection{Challenge 1: Analyze Join Timing}

Measure the time from join request to join accept:
\begin{itemize}
    \item Add timestamp logging
    \item Calculate total join time
    \item Identify RX1 vs RX2 window used
\end{itemize}

\subsection{Challenge 2: Multiple Join Attempts}

Modify the code to handle join failures:
\begin{itemize}
    \item Implement exponential backoff
    \item Retry up to 3 times
    \item Log each attempt
\end{itemize}

\subsection{Challenge 3: Display Join Parameters}

Extract and display join accept parameters:
\begin{itemize}
    \item RX1 delay
    \item RX2 data rate
    \item RX2 frequency
    \item CFList (if present)
\end{itemize}

% ============================================================================
\section{References}
% ============================================================================

\begin{itemize}
    \item LoRaWAN 1.0.4 Specification (LoRa Alliance)
    \item LoRaWAN Regional Parameters (RP002-1.0.4)
    \item Zephyr LoRaWAN API Documentation
    \item Course 1 Presentation Slides
\end{itemize}

% ============================================================================
\section{Next Steps}
% ============================================================================

After completing this lab, you should proceed to:
\begin{itemize}
    \item \textbf{Lab 1.2:} Send and receive messages
    \item \textbf{Lab 1.3:} Implement confirmed uplinks
    \item \textbf{Lab 1.4:} Downlink handling and Class A operation
\end{itemize}

\vfill

\noindent\rule{\textwidth}{0.4pt}
\begin{center}
\textit{USP Training Course 1 - Lab 1.1}\\
\textit{Version 1.0 - 2025}
\end{center}

\end{document}
